\SourceLeaf{449}%
\par\indent Differentiirt man abermals nach $s$ und bemerkt, dass
\[
\frac{\partial\partial\xi}{\partial s\partial\theta}=\frac{\partial(pX)}{\partial\theta}\qquad\text{u. s. w.}
\]
und dass
\[
X\xi'+Y\eta'+Z\zeta'=0,
\]
so wird
\[
\begin{aligned}
\frac{\partial\partial m}{\partial s^{2}}={}&p\left(\xi'\frac{\partial X}{\partial\theta}+\eta'\frac{\partial Y}{\partial\theta}+\zeta'\frac{\partial Z}{\partial\theta}\right)+p'\left(X\frac{\partial\xi}{\partial\theta}+Y\frac{\partial\eta}{\partial\theta}+Z\frac{\partial\zeta}{\partial\theta}\right)\\
={}&p\left(\xi'\frac{\partial X}{\partial\theta}+\eta'\frac{\partial Y}{\partial\theta}+\zeta'\frac{\partial Z}{\partial\theta}\right)+mp'p'\\
={}&-\left(\xi\frac{\partial X}{\partial s}+\eta\frac{\partial Y}{\partial s}+\zeta\frac{\partial Z}{\partial s}\right)\left(\xi'\frac{\partial X}{\partial\theta}+\eta'\frac{\partial Y}{\partial\theta}+\zeta'\frac{\partial Z}{\partial\theta}\right)\\
&+\left(\xi'\frac{\partial X}{\partial s}+\eta'\frac{\partial Y}{\partial s}+\zeta'\frac{\partial Z}{\partial s}\right)\left(\xi\frac{\partial X}{\partial\theta}+\eta\frac{\partial Y}{\partial\theta}+\zeta\frac{\partial Z}{\partial\theta}\right)\\
={}&\left(\frac{\partial Y}{\partial\theta}\frac{\partial Z}{\partial s}-\frac{\partial Y}{\partial s}\frac{\partial Z}{\partial\theta}\right)X+\left(\frac{\partial Z}{\partial\theta}\frac{\partial X}{\partial s}-\frac{\partial Z}{\partial s}\frac{\partial X}{\partial\theta}\right)Y\\
&+\left(\frac{\partial X}{\partial\theta}\frac{\partial Y}{\partial s}-\frac{\partial X}{\partial s}\frac{\partial Y}{\partial\theta}\right)Z.
\end{aligned}
\]

[Wird aber das zum Punkte $x,y,z$ gehörige Flächenelement
\[
mdsd\theta
\]
mittelst paralleler Normalen auf die Hülfskugel vom Radius $1$ abgebildet, so entspricht ihm daselbst ein Flächenraum von der Grösse
\[
\begin{aligned}
\Biggl\{&X\left(\frac{\partial Y}{\partial s}\frac{\partial Z}{\partial\theta}-\frac{\partial Y}{\partial\theta}\frac{\partial Z}{\partial s}\right)+Y\left(\frac{\partial Z}{\partial s}\frac{\partial X}{\partial\theta}-\frac{\partial Z}{\partial\theta}\frac{\partial X}{\partial s}\right)\\
&+Z\left(\frac{\partial X}{\partial s}\frac{\partial Y}{\partial\theta}-\frac{\partial X}{\partial\theta}\frac{\partial Y}{\partial s}\right)\Biggr\}dsd\theta;
\end{aligned}
\]
mithin ist das Krümmungsmass der Fläche in dem betrachteten Punkte gleich
\[
-\frac1m\frac{\partial\partial m}{\partial s^{2}}.\biggr]
\]
